\section{CV of the researcher}

\nb{Name:} Diederick Lambertus Vermetten\\
\nb{Date and place of Birth:} 23/08/1996, Maastricht, The Netherlands\\
\nb{Nationality:} Dutch\\
\nb{ORCID:} \href{https://orcid.org/0000-0003-3040-7162}{0000-0003-3040-7162}\\
\nb{GitHub:} \href{https://github.com/Dvermetten}{Dvermetten}

% \subsection*{Current Employment}
% I'm currently employed as a post-doctoral researcher at Sorbonne Universit\'e, CNRS, LIP6. I work on the ERC-project \emph{DynaBBO} led by Dr. Carola Doerr. 

\subsection*{Education}

\begin{itemize}
    \item[] \textit{01/01/2020–13/02/2025}: Ph.D. (Teaching PhD position), Natural Computing, Leiden University, The Netherlands. (\emph{With distinction Cum Laude}, best possible grade in the Netherlands).  \\ Thesis: \textit{From Benchmarking Optimization Heuristics to Dynamic Algorithm Configuration} \\ Advisors: Prof. Thomas Bäck, Dr. Hao Wang, Dr. Carola Doerr (external supervisor from Sorbonne Universit\'e).
    \item[] \textit{01/09/2017-16/10/2019}: MSc. in Computer Science, Leiden University, The Netherlands, Cum Laude.
    \item[] \textit{01/08/2014-31/08/2017}: BSc. Computer Science \& BSc. Mathematics, Leiden University, The Netherlands (double bachelor programme).
\end{itemize}

\subsection*{Professional Experience}
\begin{itemize}
    \item[] \textit{01/03/2025-Present}: Post-Doctoral researcher on the ERC Consolidator project `DynaBBO' at Sorbonne Universit\'e, CNRS, LIP6, France. Supervised by Carola Doerr.
    \item[] \textit{01/12/2018-30/11/2019}: Combined Teaching assistant (0.5FTE) and Software Engineer (0.5FTE) at LIACS, Leiden University, The Netherlands.
\end{itemize}

% \section{Publications}
% To copy from thesis + update


\subsection*{Supervision Activities}
At Leiden University, I have been involved in the supervision team of for \emph{5 BSc students’ theses and 10 MSc students’ theses} (including 1 external student at St. Andrews University, UK). For these projects, I have been involved in both the project conception and the daily supervision.

At LIP6, I have been working with and supervising a software engineer, Catalin Dinu. Throughout this process, I have guided him on research activities, resulting in him contributing to a published conference paper, as well as being the first author on another accepted conference paper and a paper submitted to a leading journal in evolutionary computation. 

\subsection*{Teaching Activities}
\begin{itemize}
    \item[] \textit{01/09/2019-31/06/2024}: Lecturer, Natural Computing (3rd year BSc. course). Responsible for half of the course lectures, all workgroups and assignments. 6ECTS course with 1 lecture per week during one semester.
    \item[] \textit{01/01/2020-31/06/2024}: Lecturer, Introduction to Machine Learning (1st year MSc. course). Responsible for half of the course lectures, all workgroups and assignments. 6ECTS course with 1 lecture per week during one semester. As this was a new course, I was heavily involved in the design of all course materials. 
    \item[] \textit{01/12/2018-30/11/19}: Teaching assistant for a variety of courses at the BSc and MSc level at Leiden University
\end{itemize}

\subsection*{Organization}
\begin{itemize}
    \item[] \textit{2020-Present}: Organizer of a variety of competitions at GECCO conferences (main conference of the Evolutionary Computing domain).
    \item[] \textit{2023}: Part of the online team for the EMO conference in Leiden. 
    \item[] \textit{2023}: Co-organized the SBOX-COST workshop at GECCO in 2023.
    \item[] \textit{2020}: Part of the online team of PPSN in Leiden, working on hybrid conference organization. 
    % \item[] 2025: Creator of a task-force on accessible benchmarking within the benchmarking group of COST action ROAR-NET.
\end{itemize}

\subsection*{Reviewing}
\begin{itemize}
    \item[] \textit{19/09/2024-Present}: Member of the reproducibility board of the ACM Transactions on Evolutionary Learning and Optimization (TELO)
    \item[] Reviewer for a variety of conferences and journals in the field of optimization and automated machine learning, including GECCO, CEC, TEVC, ECJ, TELO, EVOSTAR, JMLR, AutoML and more.
\end{itemize}

\subsection*{Tutorials and Invited Presentations}
\begin{itemize}
    \item[] Series of tutorials on benchmarking (focused on IOHprofiler) at GECCO, CEC
    \item[] Co-developed a tutorial series on structural bias of optimization algorithms starting at PPSN 2024
    \item[] Invited talk in JOLEA lecture series
    \item[] Lecturer at SIGEVO summer school in 2024 and 2025
    \item[] Presentations of peer-reviewed papers at GECCO, CEC, PPSN, FOGA, EVOSTAR, ECTA, AUTOML
\end{itemize}

% \todo{Paper presentations as well?}

\subsection*{Awarded Funding and Invitations}
\begin{itemize}
    \item[] \textit{2024}: Awarded funding from LUF (Leiden University Fund) for a research visit (2 weeks) to the University of Melbourne
    \item[] \textit{2022}: Awarded a \emph{SPECIES scholarship} for a 3 month stay at Edinburgh Napier University
    \item[] Invited to attend three Dagstuhl workshops: "Challenges in Benchmarking Optimization Heuristics" (23251), "Synergizing Theory and Practice of Automated Algorithm Design for Optimization" (23332), "Better Benchmarking Setups for Optimization" (25444)
    \item[] Invited to attend two Lorentz center workshops: Benchmarking in Multi-Criteria Optimisation (BeMCO), Benchmarked: Optimization Meets Machine Learning
\end{itemize}

\subsection*{Awards and Distinctions}
\begin{itemize}
    \item[] \emph{Gregorz Rozenberg award} for outstanding research during the PhD (awarded by LIACS)
    \item[] Nominated for the Leiden Science thesis award by LIACS
    \item[] 3rd place at the NeurIPS BBO-challenge 2019 (1st place on the alternative leaderboard)
    \item[] Best paper award at the ENUM track at GECCO 2023 
    \item[] Outstanding student award at EVOSTAR 2023
    \item[] 2nd place at the Niching competition in GECCO 2024 and 2025
    \item[] 2nd place at the GNBG-II competition at CEC 2025
    \item[] Best paper award at ECTA 2024 
\end{itemize}

\subsection*{Research Software Outputs}

I have been a core developer on the following projects: 
\begin{itemize}
    \item[] \href{https://github.com/IOHprofiler}{\textbf{IOHprofiler}}: a platform for analyzing and comparing iterative optimization heuristics. Combined, the different parts of IOHprofiler have achieved over $500\,000$ downloads. 
    \begin{itemize}
        \item \href{https://iohanalyzer.liacs.nl/}{IOHanalyzer}: Web-based GUI for the analysis and comparison of performance and behavior of different optimization algorithms.
        \item  \href{https://iohprofiler.github.io/IOHexperimenter/}{IOHexperimenter}: Software package (C++ and Python) which simplifies running benchmarking studies by providing robust and flexible data logging on a wide range of problems.
        \item \href{https://iohprofiler.github.io/IOHinspector/}{IOHinspector}: Python-package for local performance analysis of single- and multi-objective optimization algorithms. 
    \end{itemize}
    \item[] \href{https://github.com/nikivanstein/BIAS}{BIAS toolbox}: A Python package which integrates a wide range of statistical tests to detect structural bias in optimization heuristics.
    \item[] \href{https://github.com/Dvermetten/ModDE}{Modular Differential Evolution}: A flexible implementation of differential evolution, which enables the creation of thousands of different versions of the base algorithm. 
\end{itemize}

In addition, I have been a contributor to several projects.
\begin{itemize}
    \item[]\href{https://github.com/IOHprofiler/ModularCMAES}{Modular CMA-ES}: A flexible and efficient implementation of CMA-ES, which enables the creation of thousands of different versions of the base algorithm. This has been used in a variety of research projects. \\
    \item[]\href{https://github.com/nikivanstein/IOHxplainer}{IOHxplainer}: a toolbox which integrates methodologies from explainable machine learning into a black-box optimization context through integration with the IOHprofiler project mentioned above.\\
    \item[]\href{https://github.com/XAI-liacs/LLaMEA}{LLaMEA}: a large-language model system which integrates principles from evolutionary computation to generate and benchmark optimization algorithms, in which I integrated state-of-the-art hyperparameter optimization methodologies. 
\end{itemize}